\documentclass{article}
\usepackage[utf8]{inputenc}
\usepackage{amsmath}
\usepackage{geometry}
\geometry{margin=2cm}
\begin{document}

\section*{Questão 127 - ENEM 2024 - Ciências da Natureza}

\subsection*{Interpretação e Estratégia}

Esta questão solicita a análise do fluxograma que relaciona a urbanização a riscos para a saúde humana, especialmente por meio da transmissão de arboviroses como chikungunya e zika, ambas doenças transmitidas por mosquitos vetores.

\paragraph{Objetivo}
Identificar, a partir do fluxograma, como a urbanização contribui para o aumento dos riscos à saúde humana.

\paragraph{Estratégia}
Analisar a sequência do fluxograma e compreender a importância da redução da biodiversidade aquática para o aumento da população de mosquitos, que eleva o risco de transmissão de doenças.

\subsection*{Conteúdo e Revisão}

\begin{center}
\begin{tabular}{|p{5cm}|p{6cm}|p{6cm}|}
\hline
\textbf{Conceito} & \textbf{Chave} & \textbf{Aplicação} \\
\hline
Urbanização e saneamento básico & Urbanização sem saneamento contamina água & Relação direta no fluxograma entre urbanização e ausência de saneamento básico. \\
\hline
Biodiversidade aquática & Redução da biodiversidade diminui predadores naturais & Contaminação da água reduz predadores naturais dos mosquitos, aumentando sua população. \\
\hline
Transmissão de arboviroses & Menos predadores aumentam riscos & Maior risco epidemiológico para chikungunya e zika. \\
\hline
\end{tabular}
\end{center}

\paragraph{Fundamentos Teóricos}
\begin{itemize}
  \item Urbanização traz ausência ou deficiência de saneamento básico.
  \item Ausência de saneamento resulta em contaminação das águas urbanas.
  \item Contaminação provoca a redução da biodiversidade aquática, principalmente dos predadores naturais dos mosquitos.
  \item Redução desses predadores leva ao aumento das populações de mosquitos vetores.
  \item Populações maiores de mosquitos implicam maior risco para doenças transmitidas por eles.
\end{itemize}

\subsection*{Resolução e "Pulo do Gato"}

\begin{enumerate}
  \item O fluxograma apresenta a sequência causal: Urbanização $\to$ Ausência de saneamento básico $\to$ Contaminação da água $\to$ Redução da biodiversidade aquática $\to$ Riscos para saúde humana.
  \item A redução da biodiversidade inclui a perda dos predadores naturais dos mosquitos vetores.
  \item Menos predadores significam maior população de mosquitos.
  \item Assim, aumenta o risco de transmissão de arboviroses como chikungunya e zika.
  \item Portanto, a alternativa correta é a letra \textbf{D}.
\end{enumerate}

\paragraph{Pulo do Gato}
Conectar a ausência de saneamento à perda dos predadores naturais e ao consequente aumento dos mosquitos vetores, que elevam os riscos de doenças, esclarece o impacto da urbanização sobre a saúde pública.

\subsection*{Armadilhas na Resposta}

\begin{itemize}
  \item Confundir arboviroses com verminoses, errando o foco da transmissão.
  \item Assumir eutrofização sem respaldo no fluxograma.
  \item Valorizar presença de coliformes em detrimento da biodiversidade aquática.
  \item Focar apenas na proximidade das pessoas com mosquitos, ignorando o papel da biodiversidade aquática na regulação desses vetores.
\end{itemize}

\subsection*{Código da Questão}

CN_BIO_INTERPRETACAO_001



\end{document}